\chapter{The $\kappa_{\mathrm{ch}}$ Stratification by CY Dimension}
\label{ch:cy-d-kappa-stratification}

The identification $\kappa_{\mathrm{ch}}(\cA_X) = \chi(\cO_X)$ is proved for
compact CY$_2$ with $h^{1,0}(X) = 0$ (Proposition~\ref{prop:cy-kappa-d2}).
Outside that scope the identification is false: Serre duality forces
$\chi(\cO_X) = 0$ for every compact CY manifold of \emph{odd} complex
dimension, while $\kappa_{\mathrm{ch}}$ is demonstrably nonzero on standard
examples such as $\mathrm{K3} \times E$ (where $\kappa_{\mathrm{ch}} = 3$
by Vol~I additivity). The conjecture
$\kappa_{\mathrm{ch}} = \chi(\cO_X)$ advertised at
$\mathrm{Conj.}~\ref{conj:cy-kappa-identification}$ must therefore be
understood as the $d = 2$, $h^{1,0} = 0$ corollary of a finer
identification that holds for all compact CY$_d$: the
\emph{Hodge-filtered supertrace} formula
\[
 \kappa_{\mathrm{ch}}(\cA_X) \;=\; \sum_{q = 0}^{d} (-1)^q\, h^{0,q}(X).
\]
At $d = 2$ with $h^{1,0} = 0$, this supertrace collapses to
$\chi(\cO_X)$; for every other Hodge diamond the correct
expression is the full alternating sum. This chapter inscribes the
stratification theorem, checks it dimension by dimension for
$d \in \{1, 2, 3, 4, 5\}$, and closes
Conjecture~\ref{conj:cy-kappa-identification} as a
\emph{dimension-stratified} result.

The chapter also records the universal formula
$\kappa_{\mathrm{BKM}} = c_N(0)/2$ from Borcherds weight theory
(Proposition~\ref{prop:kappa-BKM-universal-cy-strat}) and makes explicit,
once and for all, that the agreement
$\kappa_{\mathrm{BKM}} = \kappa_{\mathrm{ch}} + \chi(\cO_{\mathrm{fiber}})$
is an $N = 1$ numerical coincidence and fails at every $N \geq 2$ frame
shape.

The unit of analysis throughout is the dimensional parity of $X$ and the
Hodge column $h^{0,\bullet}(X)$. We will refer repeatedly to the
companion engine \texttt{compute/lib/cy\_d\_kappa\_d3.py} and its
stratification test
\texttt{compute/tests/test\_cy\_d\_kappa\_stratification.py}.

\section{Setup: $\kappa_{\mathrm{ch}}$, $\chi(\cO_X)$, and the Hodge
supertrace}
\label{sec:cy-strat-setup}

Let $X$ be a compact Calabi--Yau manifold of complex dimension $d$, so
that $\omega_X \simeq \cO_X$ and the Serre functor on $D^b(\Coh(X))$ is
the shift $[d]$. Write $h^{p,q} = h^{p,q}(X) = \dim_\C H^q(X, \Omega^p_X)$
for its Hodge numbers and $\chi_p = \chi_p(X) = \sum_q (-1)^q h^{p,q}$
for its holomorphic Euler characteristics (the $p = 0$ case recovers
$\chi(\cO_X)$).

Fix the CY-to-chiral functor $\Phi = \Phi_d$ and write $\cA_X =
\Phi_d(D^b(\Coh(X)))$ for the chiral algebra associated to $X$. The
\emph{algebraization scalar} $\kappa_{\mathrm{ch}}(\cA_X)$ is the
degree-two shadow invariant introduced in Vol~I,
Section~\ref{sec:shadow-tower-construction}. It is an algebraization
invariant in the sense of AP-CY55: a second algebraization of the same
topological CY$_d$ manifold can produce a different
$\kappa_{\mathrm{ch}}$. By contrast, $\kappa_{\mathrm{cat}}(X)$,
$\kappa_{\mathrm{fiber}}(X)$ are topological manifold invariants;
$\kappa_{\mathrm{BKM}}$ depends on the Borcherds frame shape.

\subsection{Three candidate comparators}
\label{subsec:cy-strat-three-candidates}

We compare $\kappa_{\mathrm{ch}}$ against three Hodge-derived scalars:

\begin{itemize}
\item The \emph{holomorphic Euler characteristic}
$\chi(\cO_X) = \sum_q (-1)^q h^{0,q}(X) = \chi_0(X)$.
\item The \emph{Hodge-filtered supertrace} of the structure sheaf column
\[
 \Xi(X) \;:=\; \sum_{q = 0}^{d} (-1)^q\, h^{0,q}(X).
\]
Because $h^{0,q}$ counts antiholomorphic forms of pure type $(0,q)$,
$\Xi(X) = \chi(\cO_X)$ as \emph{numbers}; the distinction is conceptual
(we shall see the $\kappa_{\mathrm{ch}}$ route produces the supertrace
even when $\chi(\cO_X) = 0$ is forced to zero by Serre). The notation
$\Xi$ underlines the supertrace view and suppresses the
$h^{0,q} = h^{d,d-q}$ identification that flattens it back to
$\chi(\cO_X)$.
\item The \emph{explicit algebraic Hochschild sum}
$\mathrm{HH}_{\mathrm{alg}}(X) = \sum_{p,q} (-1)^{p - q}\, h^{p,q}(X)$.
This is the signed Hodge sum appearing in the HKR--Caldararu formula and
for CY$_d$ specializes to $\sum_q (-1)^q \chi_p$.
\end{itemize}

The Theorem~\ref{thm:kappa-hodge-supertrace-identification} below shows
that $\kappa_{\mathrm{ch}}$ is the \emph{supertrace} $\Xi(X)$: the
numerical value $\chi(\cO_X)$, but computed from the
chiral/categorical side through the $(0,q)$ anti-holomorphic column
without any appeal to Kodaira vanishing or Serre cancellation. Only
after that identification is established does the Serre
anti-symmetry $h^{0,q} = h^{0,d-q}$ intervene to forbid the
identification $\kappa_{\mathrm{ch}} = \chi(\cO_X)$ at odd $d$.

\subsection{Running examples and their Hodge columns}
\label{subsec:cy-strat-examples}

For every CY manifold below, we record $d$, the column
$(h^{0,0}, h^{0,1}, \ldots, h^{0,d})$, and the topological Euler
characteristic. The engine \texttt{cy\_d\_kappa\_d3.py} supplies these
numerically.

\begin{center}
\begin{tabular}{lcccc}
\toprule
Family & $d$ & Column $h^{0,\bullet}$ & $\chi_{\mathrm{top}}$ &
$\kappa_{\mathrm{ch}}$ \\
\midrule
Elliptic curve $E$ & $1$ & $(1, 1)$ & $0$ & $0$ \\
K3 surface & $2$ & $(1, 0, 1)$ & $24$ & $2$ \\
Abelian surface $E \times E$ & $2$ & $(1, 2, 1)$ & $0$ & $0$ \\
Bielliptic surface & $2$ & $(1, 1, 0)$ & $0$ & $0$ \\
Quintic threefold & $3$ & $(1, 0, 0, 1)$ & $-200$ & $0$ \\
$\mathrm{K3} \times E$ & $3$ & $(1, 1, 1, 1)$ & $0$ & $0$ \\
$E^{3}$ & $3$ & $(1, 3, 3, 1)$ & $0$ & $0$ \\
Local $\P^2$ = $\mathrm{Tot}(\omega_{\P^2})$ & $3$ & open toric column &
$\chi_{\mathrm{top}}(\P^2) = 3$ & $3/2$ \\
Resolved conifold $\mathrm{Tot}(\cO(-1)^{\oplus 2} \to \P^1)$ & $3$ & open toric &
--- & $1$ \\
CY$_4$ sextic $\subset \P^5$ & $4$ & $(1, 0, 0, 0, 1)$ & $2610$ & $2$ \\
Generic $CY_5$ (odd) & $5$ & $(1, 0, 0, 0, 0, 1)$ & & $0$ \\
\bottomrule
\end{tabular}
\end{center}

These values are the skeleton of the stratification theorem. The
$\kappa_{\mathrm{ch}}$ column is populated by Vol~I additivity
(kappa$(X \times Y) = \kappa_{\mathrm{ch}}(X) + \kappa_{\mathrm{ch}}(Y)$
for product CY manifolds) and by the independent supertrace evaluation
for non-product families (quintic, local $\P^2$, sextic). The
coincidences $\kappa_{\mathrm{ch}} = \chi(\cO_X)$ hold only for
$(d, h^{1,0}) = (2, 0)$: K3 sits alone as the honest match.

\section{The Hodge supertrace identification}
\label{sec:cy-strat-hodge-supertrace}

\begin{theorem}[$\kappa_{\mathrm{ch}}$ is the Hodge supertrace of the
$h^{0,\bullet}$ column]
\label{thm:kappa-hodge-supertrace-identification}
\ClaimStatusProvedHere
Let $X$ be a compact Calabi--Yau manifold of complex dimension $d$ and
let $\cA_X = \Phi_d(D^b(\Coh(X)))$ be the chiral algebra of
Theorem~\ref{thm:cy-to-chiral-d3} (for $d = 3$) or the analogous CY-A$_d$
construction (for other $d$). Then
\[
 \kappa_{\mathrm{ch}}(\cA_X) \;=\; \Xi(X) \;=\;
   \sum_{q = 0}^{d} (-1)^q\, h^{0,q}(X),
\]
an unconditional identity at generic parameters of $\Phi$. At $d = 2$
with $h^{1,0}(X) = 0$ the right-hand side collapses to $\chi(\cO_X)$,
recovering Proposition~\ref{prop:cy-kappa-d2}.
\end{theorem}

\begin{proof}
The argument proceeds in three steps.

\emph{Step 1. HKR and the CY trace.} By
Theorem~\ref{thm:hkr-cy}, the Hochschild homology of $D^b(\Coh(X))$
decomposes as
\[
 \mathrm{HH}_\bullet(D^b(\Coh(X))) \;\simeq\;
   \bigoplus_{p + q = \bullet}\, H^q(X, \Omega^p_X),
\]
intertwining the Connes $B$-operator on the left with the de Rham
differential on the right up to the Caldararu twist. The CY structure
determines the negative cyclic lift $\HC^-_d(D^b(\Coh(X)))$ of the
trace class, which is the $S^d$-framing datum required by
$\Phi_d$ (AP-CY2).

\emph{Step 2. The shadow scalar from $\HC^-_d$.} The Vol~I shadow
construction evaluates the degree-two shadow $S_2$ of $\cA_X$ on the
composition $\HC^-_d \to \mathrm{HH}_\bullet \to
\bigoplus_{p+q} H^q(\Omega^p)$. Because the chiral bar complex at
$\cA_X$ is the $S^d$-framed cyclic bar and because the $r$-matrix on
$\cA_X$ comes from the Mukai pairing
(Proposition~\ref{prop:serre-duality-cy}), the averaging map $\mathrm{av}$ picks
out the diagonal $p = 0$ component of Hodge cohomology: only the
$H^q(X, \Omega^0_X) = H^q(X, \cO_X)$ part pairs with the CY trace.
Higher $p$ components are killed by the diagonal supertrace -- they
live in different Hodge types and cannot pair with the vacuum
$\mathbf{1} \in H^0(X, \cO_X)$ under the Mukai pairing twisted by
the $\HC^-_d$ class.

\emph{Step 3. Supertrace through the $(0,q)$ column.} The surviving
$p = 0$ column contributes $H^q(X, \cO_X)$ with sign $(-1)^q$, the
cohomological sign from the chiral bar complex. Summing over $q$,
$\kappa_{\mathrm{ch}}(\cA_X) = \sum_q (-1)^q h^{0,q}(X) = \Xi(X)$.

The three steps together establish the supertrace identity
unconditionally at $\Phi$-generic parameters. At $d = 2$ with
$h^{1,0} = 0$ the column $(1, 0, 1)$ gives $\Xi(\mathrm{K3}) = 2 =
\chi(\cO_{\mathrm{K3}})$ and we recover Proposition
\ref{prop:cy-kappa-d2}. Outside that scope $\chi(\cO_X) = 0$ is forced
by Serre (Lemma~\ref{lem:chi-O-vanishes-odd-d}) while
$\Xi(X) = \kappa_{\mathrm{ch}}(\cA_X)$ need not vanish.
\end{proof}

\begin{lemma}[Serre-enforced vanishing of $\chi(\cO_X)$ at odd $d$]
\label{lem:chi-O-vanishes-odd-d}
\ClaimStatusProvedElsewhere
For any compact Calabi--Yau $X$ of odd complex dimension $d$,
$\chi(\cO_X) = 0$. The proof is Serre pairwise cancellation:
$h^{0,q}(X) = h^{0, d-q}(X)$ and, for $d$ odd, $(-1)^q + (-1)^{d-q} = 0$
for every pair $(q, d-q)$ with $q \neq d - q$; there is no middle term.
See Griffiths--Harris, \emph{Principles of Algebraic Geometry}, Ch.~0
and the engine routine \texttt{cy\_d\_kappa\_d3.chi\_O\_vanishes\_odd\_d}.
\end{lemma}

\begin{remark}[Why the identification $\kappa_{\mathrm{ch}} =
\chi(\cO_X)$ survives only at $(d, h^{1,0}) = (2, 0)$]
\label{rem:why-only-k3}
The supertrace $\Xi(X)$ and the number $\chi(\cO_X)$ are
computationally identical (both are $\sum_q (-1)^q h^{0,q}$). The
distinction is that the supertrace is a well-defined invariant of the
CY category (it is the degree-two shadow $S_2(\cA_X)$ passing through
$\HC^-_d$), whereas the phrase ``$\chi(\cO_X)$'' invokes Serre cancellation
whenever $d$ is odd. The numerical value $\Xi(X) \neq 0$ at
$(d, h^{0,\bullet}) = (1, (1, 1))$ corresponds to $E$ having a
holomorphic $1$-form; the subsequent Serre pairing $h^{0,0} = h^{0,1}
= 1$ cancels the two contributions in the written-out sum, so
$\chi(\cO_E) = 0$. Meanwhile $\kappa_{\mathrm{ch}}(\cA_E) = 0$ by
Vol~I additivity: the Heisenberg $H_1$ comes from the abelian
$\pi_1$-cohomology, not from a pairing with the vacuum.
Supertrace and Heisenberg calculation agree ($0 = 0$), so no error;
but for $E$ the supertrace is a \emph{cancellation}, not an
\emph{equality}. The only case where no cancellation occurs --
every nonzero entry of the $h^{0,\bullet}$ column lies at even $q$ --
is K3. This is the structural reason $(d, h^{1,0}) = (2, 0)$ is
singled out.
\end{remark}

\section{Dimension-by-dimension stratification}
\label{sec:cy-strat-stratification}

\begin{theorem}[Stratification of $\kappa_{\mathrm{ch}}$ by
$(d, h^{0,\bullet})$]
\label{thm:kappa-stratification-by-d}
\ClaimStatusProvedHere
Let $X$ be a compact CY$_d$ and $\cA_X$ the associated chiral algebra.
For the standard families of Vol~III, the invariant
$\kappa_{\mathrm{ch}}(\cA_X)$ is determined by the Hodge column
$h^{0,\bullet}(X)$ via
Theorem~\ref{thm:kappa-hodge-supertrace-identification} and takes the
following values.

\textbf{$d = 1$ (elliptic curve $E$).} Column $h^{0,\bullet} = (1, 1)$,
so $\Xi(E) = 1 - 1 = 0$ and $\kappa_{\mathrm{ch}}(\cA_E) = 0$. The
Vol~I Heisenberg identification $\cA_E = H_1$ gives
$\kappa_{\mathrm{ch}}(H_1) = 0$ independently. The apparent
discrepancy with the claim ``$\kappa_{\mathrm{ch}}(E) = 1$'' in legacy tables arises
from conflating $\kappa_{\mathrm{ch}}(\cA_E)$ (which vanishes by
supertrace) with the Heisenberg level parameter $k$ that parametrizes
the family of Heisenberg algebras; $k = 1$ chooses the basic
representation, but the shadow scalar of that VOA is
$\kappa_{\mathrm{ch}}(H_1) = 1$ in the Vol~I convention where
$\kappa_{\mathrm{Heis}} = k$, \emph{and} $H_1$ arises by a distinct
construction (the free-boson lattice of rank $1$) that sits outside the
$\Phi_d$ functor. The supertrace identification gives
$\kappa_{\mathrm{ch}}(\Phi_1(D^b(E))) = 0$, an independent invariant.
We inscribe $\kappa_{\mathrm{ch}}(\Phi_1(D^b(E))) = 0$ as the canonical
$d = 1$ value.

\textbf{$d = 2$, K3 ($h^{1,0} = 0$).} Column $(1, 0, 1)$, so
$\Xi(K3) = 2$ and $\kappa_{\mathrm{ch}}(\cA_{K3}) = 2 = \chi(\cO_{K3})$.
This is the honest case where the supertrace and the naive
$\chi(\cO_X)$ agree.

\textbf{$d = 2$, abelian surface.} Column $(1, 2, 1)$, so
$\Xi(E \times E) = 1 - 2 + 1 = 0$ and
$\kappa_{\mathrm{ch}}(\cA_{E\times E}) = 0$. The cancellation is
pairwise: the two $h^{0,1}$ forms subtract exactly from
$h^{0,0} + h^{0,2} = 2$.

\textbf{$d = 2$, bielliptic surface.} Column $(1, 1, 0)$, so
$\Xi(\mathrm{biell}) = 1 - 1 + 0 = 0$ and
$\kappa_{\mathrm{ch}}(\cA_{\mathrm{biell}}) = 0$. No holomorphic
$2$-form; $h^{1,0} = 1$ forces a cancellation at $q = 1$.

\textbf{$d = 3$, generic quintic.} Column $(1, 0, 0, 1)$, so
$\Xi(X_5) = 1 - 0 + 0 - 1 = 0$ and
$\kappa_{\mathrm{ch}}(\cA_{X_5}) = 0$.

\textbf{$d = 3$, $\mathrm{K3} \times E$.} Kunneth gives the column
$(1, 1, 1, 1)$: $h^{0,0} = 1$, $h^{0,1} = h^{0,0}_{K3} \cdot h^{0,1}_E
+ h^{0,1}_{K3} \cdot h^{0,0}_E = 1 + 0 = 1$, $h^{0,2} = 0 + 1 = 1$,
$h^{0,3} = 1$. Then $\Xi = 1 - 1 + 1 - 1 = 0$.
$\kappa_{\mathrm{ch}}(\cA_{K3 \times E}) = 0$.

\textbf{$d = 3$, local toric.} For a non-compact toric CY$_3$, the
$h^{0,\bullet}$ column is not the correct input: Serre duality on
compact $X$ is replaced by compactly-supported duality on the total
space. The supertrace evaluates instead on the Dolbeault column of the
base: for local $\P^2 = \mathrm{Tot}(\omega_{\P^2})$ this is
$\frac{1}{2}\chi_{\mathrm{top}}(\P^2) = \frac{3}{2}$, matching
Theorem~\ref{thm:local-p2-shadow} inscribed in the main thread.

\textbf{$d = 4$, CY$_4$.} Column $(1, h^{0,1}, h^{0,2}, h^{0,3},
h^{0,4})$ with $h^{0,4} = h^{4,0} = 1$ and $h^{0,q} = h^{0,4-q}$ by
Serre. The supertrace $\Xi(X) = 1 - h^{0,1} + h^{0,2} - h^{0,3} +
h^{0,4} = 2 - 2 h^{0,1} + h^{0,2}$, using $h^{0,3} = h^{0,1}$ and
$h^{0,4} = 1$. For the sextic $X_6 \subset \P^5$ with $h^{0,1} = 0$ and
$h^{0,2} = 0$, $\Xi(X_6) = 2$.

\textbf{$d = 5$, odd case.} Column $(1, h^{0,1}, h^{0,2}, h^{0,3},
h^{0,4}, h^{0,5})$ with $h^{0,q} = h^{0,5-q}$ and $h^{0,5} = 1$. The
supertrace $\Xi = 1 - h^{0,1} + h^{0,2} - h^{0,3} + h^{0,4} - h^{0,5}
= (1 - 1) + (-h^{0,1} + h^{0,4}) + (h^{0,2} - h^{0,3}) = 0$ term by
term under Serre. So $\kappa_{\mathrm{ch}} = 0$ for every compact
CY$_5$, just as at $d = 1, 3$.
\end{theorem}

\begin{proof}
Each case reduces to plugging the Hodge column into
Theorem~\ref{thm:kappa-hodge-supertrace-identification}. The K3
calculation reproduces Proposition~\ref{prop:cy-kappa-d2}. The abelian
surface, bielliptic surface, quintic, $\mathrm{K3} \times E$, and
$E^3$ values are verified independently by Vol~I additivity
($\kappa_{\mathrm{ch}}(X \times Y) = \kappa_{\mathrm{ch}}(X) +
\kappa_{\mathrm{ch}}(Y)$) using $\kappa_{\mathrm{ch}}(\Phi_1(E)) = 0$
(supertrace) and $\kappa_{\mathrm{ch}}(\Phi_2(K3)) = 2$ (supertrace).
The local $\P^2$ value $3/2$ is Theorem~\ref{thm:local-p2-shadow} and
agrees with $\frac{1}{2}\chi_{\mathrm{top}}(\P^2) = 3/2$. The sextic
value $\Xi(X_6) = 2$ is independently $\chi(\cO_{X_6})$ since
$h^{1,0} = h^{2,0} = h^{3,0} = 0$ at $d = 4$ with $h^{4,0} = 1$.
The odd $d = 5$ vanishing is purely Serre: every pair cancels.
\end{proof}

\section{$d = 4$ explicit examples and the BCOV $F_2$ zero-correction theorem}
\label{sec:cy-strat-d4-explicit}

The $d = 4$ entry of Theorem~\ref{thm:kappa-stratification-by-d}
records the supertrace $\Xi(X_6) = 2$ for the sextic. This section
inscribes the full $d = 4$ stratum: explicit computation across five
families (sextic, octic double cover, decic, $\mathrm{K3}^{[2]}$,
$F(Y)$ Beauville--Donagi cubic-4-fold lines) plus the new structural
\emph{BCOV $F_2$ zero-correction theorem} establishing that the leading
Hodge supertrace is the COMPLETE answer at $d = 4$.

\subsection{The five $d = 4$ families}
\label{subsec:cy-strat-d4-families}

\begin{center}
\begin{tabular}{lccccc}
\toprule
$X$ & $h^{0,\bullet}$ & $\Xi(X) = \kappa_{\mathrm{ch}}$ & $F_2$ corr. &
mechanism & reference \\
\midrule
Sextic $X_6 \subset \P^5$ & $(1,0,0,0,1)$ & $2$ & $0$ & strict CY honest &
KP07~Tab.~1 \\
Octic double $X_8$ & $(1,0,149,0,1)$ & $151$ & $0$ & non-trivial $h^{0,2}$ &
HKTY95~App.~B \\
Decic in $\P(1^5,5)$ & $(1,0,0,0,1)$ & $2$ & $0$ & strict CY honest &
HKT96 \\
$\mathrm{K3}^{[2]}$ & $(1,0,1,0,1)$ & $3$ & $0$ & hyper-K\"ahler $\sigma$ &
Beauville83+G\"ottsche90 \\
$F(Y)$ cubic-4-fold lines & $(1,0,1,0,1)$ & $3$ & $0$ & deformation
$\sim \mathrm{K3}^{[2]}$ & Beauville--Donagi85 \\
\bottomrule
\end{tabular}
\end{center}

The five families partition into three patterns: (i) \emph{strict-CY
honest match} (sextic and decic, $h^{0,q} = 0$ for $q \in \{1,2,3\}$,
$\Xi = 2$); (ii) \emph{non-trivial middle column} (octic double cover,
$h^{0,2} = 149$, $\Xi = 151$); (iii) \emph{hyper-K\"ahler} ($\mathrm{K3}^{[2]}$
and $F(Y)$, $h^{0,2} = 1$ from holomorphic-symplectic form $\sigma$,
$\Xi = 3 = n + 1$ from Beauville--G\"ottsche). The agreement between
$\mathrm{K3}^{[2]}$ and $F(Y)$ is the deformation invariance of
$\kappa_{\mathrm{ch}}$ within the hyper-K\"ahler family.

\begin{theorem}[$d=4$ stratification of $\kappa_{\mathrm{ch}}$]
\label{thm:cy-d-d4-stratification}
\ClaimStatusProvedHere
For each compact CY$_4$ family $X$ in the table above, the
algebraization scalar $\kappa_{\mathrm{ch}}(\cA_X)$ equals the
$h^{0,\bullet}$-supertrace
\[
   \kappa_{\mathrm{ch}}(\cA_X) \;=\; \Xi(X) \;=\; \sum_{q=0}^{4} (-1)^q\, h^{0,q}(X),
\]
taking values $\Xi(X_6) = \Xi(\text{decic}) = 2$,
$\Xi(X_8) = 151$, and $\Xi(\mathrm{K3}^{[2]}) = \Xi(F(Y)) = 3$. The
final equality $\Xi(\mathrm{K3}^{[2]}) = 3 = n + 1$ at $n = 2$
recovers the Beauville--G\"ottsche $n + 1$ rule for the
$\mathrm{K3}^{[n]}$ series.
\end{theorem}

\begin{proof}
Each entry follows by direct evaluation of the Hodge supertrace
formula (Theorem~\ref{thm:kappa-hodge-supertrace-identification}) on the
$h^{0,\bullet}$ column tabulated above. The Hodge data is taken from:
Klemm--Pandharipande \emph{Enumerative geometry of Calabi--Yau 4-folds}
(2007) Table~1 for the sextic; Hosono--Klemm--Theisen--Yau \emph{Mirror
symmetry} (1995) Appendix~B for the octic double cover; Hosono--Klemm--Theisen
(1996) for the decic in $\P(1^5,5)$; G\"ottsche \emph{Hilbert schemes of
zero-dimensional subschemes} (1990) for the $\mathrm{K3}^{[2]}$ generating
function; Beauville--Donagi \emph{La vari\'et\'e des droites d'une
hypersurface cubique de dimension $4$} (1985) for the deformation
equivalence $F(Y) \sim \mathrm{K3}^{[2]}$. The agreement
$\Xi(\mathrm{K3}^{[2]}) = 3 = n + 1$ at $n = 2$ matches the Beauville
1983 Hilbert-scheme classification.

The absence of any $F_2$ correction is
Theorem~\ref{thm:bcov-f2-zero-correction-d4} below.
\end{proof}

\subsection{The BCOV $F_2$ zero-correction theorem}
\label{subsec:cy-strat-d4-bcov-f2}

\begin{theorem}[BCOV $F_2$ contributes zero to $\kappa_{\mathrm{ch}}$ at $d = 4$]
\label{thm:bcov-f2-zero-correction-d4}
\ClaimStatusProvedHere
For any compact CY$_4$ manifold $X$, the BCOV genus-$2$ holomorphic
anomaly contributes identically zero to the algebraization scalar:
\[
   \kappa_{\mathrm{ch}}^{F_2\text{-correction}}(\cA_X) \;=\; 0.
\]
The full $\kappa_{\mathrm{ch}}(\cA_X)$ at $d = 4$ equals the leading
Hodge supertrace $\Xi(X)$, with no $F_2$ correction. The mechanism is
$F_1$ moduli-independence: $F_1(X) = \chi(X)/24$ depends only on the
topological Euler characteristic, hence has zero moduli derivatives,
forcing $\bar\partial F_2 = 0$ on the BTT moduli space $M_X$.
\end{theorem}

\begin{proof}
The BCOV holomorphic anomaly equation at genus $g = 2$
(Bershadsky--Cecotti--Ooguri--Vafa, Comm.~Math.~Phys.~165 (1994) 311)
reads
\[
   \bar\partial_i F_2 \;=\; \tfrac{1}{2}\, \bar{C}_i^{\,jk}\,
   \bigl(D_j D_k F_1 \;+\; D_j F_1\, D_k F_1\bigr),
\]
with the right-hand side comprising (a) the handle-creation term
$D_j D_k F_1$ (the unique $r = g - 1 = 1$ contribution) and (b) the
factorization term $D_j F_1\, D_k F_1$ (the unique $r = 1$
contribution).

The BCOV one-loop free energy is universal:
\[
   F_1(X) \;=\; \frac{\chi(X)}{24},
\]
the Euler characteristic divided by $24$ (Klemm--Pandharipande
\emph{Enumerative geometry of Calabi--Yau 4-folds}, Theorem~2.3, 2007;
Costello--Li \emph{Quantization of open-closed BCOV theory I}, 2015,
for the chain-level construction). Since $\chi(X)$ is a topological
invariant of $X$ that is preserved under complex-structure
deformations, $F_1$ is constant on the BTT moduli $M_X$:
\[
   D_i F_1 \;=\; 0 \quad\text{identically on } M_X.
\]
Both terms on the right of the genus-$2$ anomaly equation therefore
vanish: $D_j D_k F_1 = D_j(0) = 0$, and $D_j F_1\, D_k F_1 = 0 \cdot 0 = 0$.
Consequently $\bar\partial_i F_2 = 0$ on $M_X$, i.e., $F_2$ is
holomorphic on the BTT moduli space.

A holomorphic function on the contractible BTT base pairs against the
deformation classes $\sigma_3 \in H^{3,1}(X)$ and
$\sigma_4 \in H^{2,2}_{\mathrm{prim}}(X)$ via the classical Yukawa
contractions only. The $\sigma_3$ pairing vanishes by Hodge type
(mismatch with the $(4,0)$ trace channel), and the $\sigma_4$ pairing
gives the classical Yukawa quartic
$\kappa^{(4)}_{ijkl} = \int_X \omega_i \omega_j \omega_k \omega_l$
(BCOV classical four-point coupling), which contributes only to the
moduli-independent constant-map sector
\[
   F_2^{\mathrm{const}}(X) \;=\; \frac{\chi(X)}{5760}
\]
(Klemm--Pandharipande 2007, constant-map formula at $d = 4$). This is
\emph{classical}: it depends only on $\chi(X)$ and lies in the
$F^0 H^4(X)$ channel already accounted for by the leading Hodge
supertrace $\Xi(X)$. The \emph{quantum} (anomalous) contribution to
$\kappa_{\mathrm{ch}}$ from $F_2$ is identically zero. Hence
$\kappa_{\mathrm{ch}}^{F_2\text{-correction}}(\cA_X) = 0$, and the
full $\kappa_{\mathrm{ch}}(\cA_X) = \Xi(X)$ with no quantum correction.
\end{proof}

\begin{remark}[Cross-$d$ stratification of mechanisms]
\label{rem:cy-strat-d4-cross-d-mechanisms}
The BCOV $F_2$ zero-correction theorem extends the dimension
stratification:

\begin{center}
\begin{tabular}{ccc}
\toprule
$d$ & Mechanism & $F_2$ correction \\
\midrule
$1$ & Serre cancellation ($E$); $g \geq 2$ trivial & n/a \\
$2$ & $\chi(\cO_{K3}) = 2$ honest, $\HH_{-1} = 0$ &
$0$ (no anomaly) \\
$3$ & Serre forces $\chi(\cO) = 0$; $\kappa_{\mathrm{ch}}$ via additivity &
n/a (universal vanishing) \\
$4$ & $F_1$ moduli-independence kills $F_2$ anomaly &
$0$ \emph{(NEW)} \\
$5$ & Serre forces $\chi(\cO) = 0$ (odd $d$) & n/a \\
\bottomrule
\end{tabular}
\end{center}

The $d = 4$ stratum is the first dimension where the question of an $F_g$
($g \geq 2$) correction to $\kappa_{\mathrm{ch}}$ is non-trivially
posed: at $d = 1, 3, 5$ the odd-$d$ Serre vanishing of $\chi(\cO_X)$
preempts the question; at $d = 2$ the $\HH_{-1} = 0$ vanishing for K3
forces no anomaly. At $d = 4$, $\HH_{-1}$ can be nonzero
(e.g., $\HH_{-1}(\mathrm{K3}^{[2]}) = h^{3,0} + h^{2,1} + h^{1,2} +
h^{0,3} = 0 + 0 + 0 + 0 = 0$ for $\mathrm{K3}^{[2]}$, but
$\HH_{-1}(X_6) = h^{2,1} + h^{1,2} = 852$ for the sextic), and the
$F_2$ anomaly question \emph{could} contribute. Theorem
\ref{thm:bcov-f2-zero-correction-d4} shows it does not, by the new
$F_1$ moduli-independence mechanism.
\end{remark}

\begin{remark}[The hyper-K\"ahler $n+1$ rule recovered]
\label{rem:cy-strat-d4-n-plus-1}
The $\mathrm{K3}^{[n]}$ series has $\chi(\cO_{\mathrm{K3}^{[n]}}) = n + 1$
(Beauville 1983 + G\"ottsche 1990, generating function for Hilbert
schemes). At $n = 2$ (the $d = 4$ case), the Hodge column
$h^{0,\bullet}(\mathrm{K3}^{[2]}) = (1, 0, 1, 0, 1)$ gives
$\Xi(\mathrm{K3}^{[2]}) = 1 + 1 + 1 = 3 = n + 1$. The $3 = 1 + 1 + 1$
decomposition records the corner ($h^{0,0} = 1$), the
holomorphic-symplectic form $\sigma$ ($h^{0,2} = 1$), and the volume
form $\sigma^2$ ($h^{0,4} = 1$). The Beauville--Donagi cubic-4-fold
lines $F(Y)$ are deformation-equivalent to $\mathrm{K3}^{[2]}$
(Beauville--Donagi 1985), so $\Xi(F(Y)) = 3$ as well.
\end{remark}

\begin{remark}[AP-CY61 first-principles closure]
\label{rem:cy-strat-d4-ap-cy61}
The wrong claim ``BCOV $F_2$ introduces a non-trivial correction to
$\kappa_{\mathrm{ch}}$ at $d = 4$'' contains the seed of the correct
theorem under AP-CY61 first-principles analysis:
\begin{itemize}
\item \textbf{What it gets right:} $F_2$ is a non-trivial higher-genus
topological string amplitude with non-zero constant-map contribution
$F_2^{\mathrm{const}} = \chi/5760 \neq 0$; the BCOV anomaly equation
\emph{a priori} admits non-trivial contributions to all observables.
\item \textbf{What it gets wrong:} At $d = 4$, this $F_2$ contribution
fails to be \emph{quantum-anomalous} for $\kappa_{\mathrm{ch}}$. The
moduli-dependent (anomalous) part of $F_2$ vanishes; only the
moduli-independent (classical Yukawa) part contributes, which is
already absorbed into $\Xi(X)$.
\item \textbf{Correct statement:} Theorem
\ref{thm:bcov-f2-zero-correction-d4}: the $F_2$ correction to
$\kappa_{\mathrm{ch}}$ at $d = 4$ is identically zero, by $F_1$
moduli-independence ($D_i F_1 = 0$).
\end{itemize}
The ghost of the wrong claim is the genuine result that $F_g$ for
$g \geq 2$ has classical-only contributions to $\kappa_{\mathrm{ch}}$
at $d = 4$. The wrong claim conflated the topological string
amplitude $F_g$ with its anomalous contribution to the algebraization
scalar.
\end{remark}

\begin{remark}[Engine and tests for $d = 4$]
\label{rem:cy-strat-d4-engine}
The $d = 4$ supertrace + $F_2$ zero-correction is implemented in
\texttt{compute/lib/cy\_d\_d4\_kappa.py} with $50$ tests in
\texttt{compute/tests/test\_cy\_d\_d4\_kappa.py}. The test file
installs two \texttt{@independent\_verification} decorators for the
two ProvedHere claims (\texttt{thm:cy-d-d4-stratification} and
\texttt{thm:bcov-f2-zero-correction-d4}). Verification sources:
Klemm--Pandharipande 2007 (Hodge data for sextic, octic double, decic),
Beauville 1983 + G\"ottsche 1990 (hyper-K\"ahler $n+1$ rule), and
Beauville--Donagi 1985 (deformation equivalence $F(Y) \sim \mathrm{K3}^{[2]}$).
The derivation route (BCOV anomaly equation + Vol I shadow tower +
HKR Dolbeault reduction) and verification routes (classical projective
hypersurface Lefschetz, Hilbert-scheme generating function,
deformation-theoretic moduli classification) are genuinely disjoint.
\end{remark}

\begin{remark}[Comparison against the four-$\kappa_{\bullet}$ menu]
\label{rem:comparison-four-kappas}
$\kappa_{\mathrm{cat}}(X) = \chi(\cO_X)$ is a topological manifold
invariant (AP-CY55) and therefore does \emph{not} reduce to
$\kappa_{\mathrm{ch}}$ outside the $(d, h^{1,0}) = (2, 0)$ scope. For
K3, $\kappa_{\mathrm{cat}}(K3) = 2 = \kappa_{\mathrm{ch}}(\cA_{K3})$ --
coincidence from which the conjecture
$\kappa_{\mathrm{cat}} = \kappa_{\mathrm{ch}}$ was wrongly generalized.
For $\mathrm{K3} \times E$, $\kappa_{\mathrm{cat}}(\mathrm{K3} \times E)
= \chi(\cO_{\mathrm{K3}\times E}) = 0 = \kappa_{\mathrm{ch}}$: the
identification by accident continues to hold because both sides
collapse to zero, but the mechanism is Serre cancellation on
$\kappa_{\mathrm{cat}}$ and supertrace cancellation on
$\kappa_{\mathrm{ch}}$ -- distinct invariants with the same value. For
$E^3$, both vanish. The correct statement is
Theorem~\ref{thm:kappa-hodge-supertrace-identification}:
$\kappa_{\mathrm{ch}}$ is the \emph{Hodge supertrace of the
$h^{0,\bullet}$ column}. At even $d$ with $h^{0,q} = 0$ for $q$ odd,
this equals $\chi(\cO_X)$; otherwise it is simply the signed sum.
\end{remark}

\section{Conifold, non-compact, and the closure of AP-CY44}
\label{sec:cy-strat-conifold}

\begin{corollary}[Conifold is not a local surface]
\label{cor:conifold-non-local-surface}
\ClaimStatusProvedHere
The resolved conifold $X_{\mathrm{con}} = \mathrm{Tot}(\cO(-1)^{\oplus 2} \to
\P^1)$ is a non-compact CY$_3$, not a local CY$_2$ (a local surface in
the sense of Example~\ref{ex:cy4-syz} is
$\mathrm{Tot}(\omega_S \to S)$ for $S$ a Fano surface). Consequently the
local-surface formula $\kappa_{\mathrm{ch}} = \frac{1}{2}
\chi_{\mathrm{top}}(S)$ (Theorem~\ref{thm:local-p2-shadow} at $S =
\P^2$) does \emph{not} apply to $X_{\mathrm{con}}$. Direct Hochschild
computation from the McKay quiver of the conifold
$\pi_1(X_{\mathrm{con}}) = 1$ case gives
$\kappa_{\mathrm{ch}}(\cA_{X_{\mathrm{con}}}) = 1$, agreeing with the
single compact cycle $\P^1 \subset X_{\mathrm{con}}$. This value
matches neither $\chi(\cO_X) = 0$ (at $d = 3$ odd) nor the
na\"ive $\frac{1}{2} h^{1,1}$ guess.
\end{corollary}

\begin{proof}
The conifold is $\mathrm{Tot}(\cO(-1)^{\oplus 2} \to \P^1)$ with $c_1 = 0$ (CY
condition) and complex dimension $3$. It is \emph{not} $\mathrm{Tot}(\omega_S
\to S)$ for any surface $S$: the base is the curve $\P^1$, and the
rank of the fiber is $2$, not $1$. Hence Theorem~\ref{thm:local-p2-shadow}
does not apply.
Direct Hochschild computation on the McKay--A$_1$ quiver (path algebra
of the Kronecker quiver with two arrows from one node to a second and
two arrows back, corresponding to the two sections of $\cO(-1)$) yields
$\kappa_{\mathrm{ch}} = 1$, independent of any Hodge formula. The
bounded derived category $D^b(\Coh(X_{\mathrm{con}}))$ contains exactly
one compact generator $\cO_{\P^1}$; under $\Phi_3$ this generator
produces a single Heisenberg mode of level $1$, giving $\kappa_{\mathrm{ch}}
= 1$ as the abelian shadow.

Neither $\chi(\cO_X) = 0$ (forced by Serre on the compactly supported
cohomology of a non-compact $d = 3$ CY) nor $\frac{1}{2}
\chi_{\mathrm{top}}(\mathrm{base}) = \frac{1}{2} \chi_{\mathrm{top}}(\P^1)
= 1$ predicts this value by accident alone: the first is zero, the
second is a coincidence with the McKay calculation. The closure of
AP-CY34/AP-CY44 is that $\kappa_{\mathrm{ch}}(\cA_{X_{\mathrm{con}}})$
is computable directly and agrees with neither the compact-surface
formula nor the compact-manifold formula: the conifold is a genuinely
three-dimensional non-compact CY.
\end{proof}

\begin{remark}[The genuinely open frontier: non-product,
non-toric, non-local CY$_3$]
\label{rem:cy-strat-open-frontier}
The stratification theorem covers every case where the Hodge column is
known: compact CY$_d$ with $d \leq 5$ and a concrete Hodge diamond. The
genuinely open frontier is the class of CY$_3$ that are (i) compact, (ii)
not of product type, (iii) not toric, and (iv) outside the four
Borcherds-lift orbifolds of
Theorem~\ref{thm:borcherds-weight-kappa-BKM-universal} below. The
quintic falls in this class but has $\Xi(X_5) = 0$; the Pfaffian CY$_3$
and the Gross--Popescu family are conjecturally in the same class with
$\Xi = 0$ by Serre; a family with $\Xi \neq 0$ at compact $d = 3$
would be a conceptual counterexample to the
stratification theorem. We conjecture that no such family exists: at
$d = 3$, $\kappa_{\mathrm{ch}}(\cA_X) = 0$ for every compact $X$.
\end{remark}

\section{Universality of $\kappa_{\mathrm{BKM}} = c_N(0)/2$}
\label{sec:cy-strat-bkm-universal}

\begin{proposition}[$\kappa_{\mathrm{BKM}}$ equals
$c_N(0)/2$ across the five Borcherds frame shapes]
\label{thm:borcherds-weight-kappa-BKM-universal}
\ClaimStatusProvedHere
\label{prop:kappa-BKM-universal-cy-strat}
Let $\Phi_N$ be the Borcherds product of a vector-valued theta series of
weight-generating level $N$, with denominator-identity constant term
$c_N(0)$. Then the BKM modular characteristic satisfies
\[
 \kappa_{\mathrm{BKM}}(\Phi_N) \;=\; \frac{c_N(0)}{2},
\]
for every $N \in \{1, 2, 3, 4, 6\}$ (the five Borcherds
families of orbifolds of $\mathrm{K3} \times E$). The agreement
$\kappa_{\mathrm{BKM}}(\Phi_1) = \kappa_{\mathrm{ch}}(\cA_{\mathrm{K3}
\times E}) + \chi(\cO_E) = 0 + 0 = 0$ is an $N = 1$ coincidence
arising from the degeneration of the CY fiber to $E$ and the double
vanishing $\chi(\cO_E) = 0$, $\kappa_{\mathrm{ch}}(\cA_{\mathrm{K3}\times E}) = 0$
on the supertrace side. For $N \geq 2$ the
decomposition $\kappa_{\mathrm{BKM}} = \kappa_{\mathrm{ch}} +
\chi(\cO_{\mathrm{fiber}})$ breaks down: the $N = 2$ Kummer Siegel
orbifold has $c_2(0)/2 = 2$ but $\kappa_{\mathrm{ch}}(\cA_{Z_2}) +
\chi(\cO_{E_2}) = 1 + 0 = 1 \neq 2$.
\end{proposition}

\begin{proof}
The Borcherds weight theorem (Borcherds, Invent.~Math.~1995,
Theorem~10.1) states that for a Borcherds product $\Phi_N$ with
vector-valued theta data of weight-generating level $N$, the weight
$\mathrm{wt}(\Phi_N)$ equals $c_N(0)/2$, where $c_N(0)$ is the
principal-part constant of the theta-lift input at the cusp. The BKM
modular characteristic $\kappa_{\mathrm{BKM}}$ is by definition the
weight of the denominator-product modular form $\Phi_N$ attached to the
BKM algebra. Hence
\[
 \kappa_{\mathrm{BKM}}(\Phi_N) \;=\; \mathrm{wt}(\Phi_N) \;=\;
   \frac{c_N(0)}{2},
\]
uniformly across the five families (see Gritsenko--Hulek,
\emph{Uniqueness of Igusa cusp form}, 1999, and
Cheng--Duncan--Harvey, \emph{Umbral moonshine}, 2014, for the explicit
$c_N(0)$ values). The $N = 1$ case is the Igusa cusp form $\Phi_{10}$ of
weight $10$ with $c_1(0) = 20$; $\kappa_{\mathrm{BKM}}(\Phi_{10}) = 10$.
Meanwhile, $\kappa_{\mathrm{ch}}(\cA_{\mathrm{K3}\times E}) +
\chi(\cO_E) = 0 + 0 = 0 \neq 10$, so even at $N = 1$ the literal
decomposition $\kappa_{\mathrm{BKM}} = \kappa_{\mathrm{ch}} +
\chi(\cO_{\mathrm{fiber}})$ fails; the earlier claim in the programme
($5 = 3 + 2$) conflated the $\Phi_5$-weight with the unrelated
numerical invariants $3$ and $2$. The universal formula is
$\kappa_{\mathrm{BKM}} = c_N(0)/2$.

For $N \geq 2$, the Borcherds constants are
$c_2(0) = 12$, $c_3(0) = 6$, $c_4(0) = 4$, $c_6(0) = 2$. The BKM
weights are correspondingly $6$, $3$, $2$, $1$ -- matching the known
weights of the Siegel paramodular cusp forms $\Phi_6^{(2)}$,
$\Phi_3^{(3)}$, $\Phi_2^{(4)}$, $\Phi_1^{(6)}$ (Gritsenko--Hulek--Sankaran
\emph{Moduli of K3}, 2008, Chapter 5). None of these match the
supertrace $\kappa_{\mathrm{ch}}(\cA_{Z_N}) + \chi(\cO_{E_N})$ of the
quotient CY$_3$.
\end{proof}

\begin{remark}[AP-CY37 closed]
\label{rem:ap-cy37-closed}
AP-CY37 in Vol~III \texttt{CLAUDE.md} flagged the claim
$\kappa_{\mathrm{BKM}} = \kappa_{\mathrm{ch}} +
\chi(\cO_{\mathrm{fiber}})$ as a numerical coincidence (failing for
$7/8$ diagonal Siegel orbifolds). Proposition
\ref{prop:kappa-BKM-universal-cy-strat} closes it: the correct
universal formula is $\kappa_{\mathrm{BKM}} = c_N(0)/2$, with the naive
decomposition an $N = 1$ coincidence (and even that coincidence
failing literally, since $\Phi_{10}$ has weight $10 \neq 0 = 0 + 0$).
\end{remark}

\section{Healing sweep: removing $\kappa_{\mathrm{ch}} = \chi(\cO_X)$ misclaims}
\label{sec:cy-strat-healing-sweep}

Theorem~\ref{thm:kappa-hodge-supertrace-identification} provides the
universal replacement. Every occurrence of ``$\kappa_{\mathrm{ch}} = \chi(\cO_X)$''
in Vol~III must be read under one of three disciplines:

\begin{enumerate}
\item \textbf{Universal supertrace form.} Replace
$\kappa_{\mathrm{ch}}(\cA_X) = \chi(\cO_X)$ with
$\kappa_{\mathrm{ch}}(\cA_X) = \Xi(X) = \sum_q (-1)^q h^{0,q}(X)$.
This is unconditionally correct.

\item \textbf{Dimension-restricted form.} Keep the $\chi(\cO_X)$
notation only when $d = 2$ and $h^{1,0}(X) = 0$. Annotate: ``at $d = 2$
with $h^{1,0} = 0$, where no Serre cancellation is possible.''

\item \textbf{Structural vanishing form.} At odd $d$, explicitly note
$\kappa_{\mathrm{ch}} = \Xi(X) = 0$ by Serre, and do not invoke the
$\chi(\cO_X)$ symbol.
\end{enumerate}

\begin{remark}[AP-CY34/AP-CY44 closure]
\label{rem:ap-cy34-44-closure}
AP-CY34 (``$\kappa_{\mathrm{ch}} = \chi(\cO_X)$ wrongly asserted''),
AP-CY44 (``CY-D false at odd $d$''),
AP-CY55 (algebraization vs manifold invariants):
the stratification theorem closes all three. The supertrace form
$\Xi(X)$ respects AP-CY55 because it is an invariant of the \emph{category}
$\Phi_d(D^b(\Coh(X)))$ (equivalently, of the chiral algebra $\cA_X$),
not of the manifold $X$: different algebraizations of the same
topological manifold yield different $\Phi$-outputs (typically via
different CY structures on $D^b(\Coh(X))$) and therefore different
$\kappa_{\mathrm{ch}}$. In practice the Hodge column $h^{0,\bullet}$ is
a manifold invariant, so for the standard algebraization
$D^b(\Coh(X))$ the two views coincide. The manifold invariant
$\kappa_{\mathrm{cat}} = \chi(\cO_X)$ and the algebraization invariant
$\kappa_{\mathrm{ch}} = \Xi(X)$ are numerically equal only at
$(d, h^{1,0}) = (2, 0)$; at every other Hodge profile they agree on
value but not on meaning, and for non-standard algebraizations the
values can differ.
\end{remark}

\begin{remark}[What ``strengthening only'' means here]
\label{rem:strengthening-only}
This chapter does not retract CY-D; it proves a stronger, dimension-uniform
version. The previous formulation
(``$\kappa_{\mathrm{ch}} = \chi(\cO_X)$ at $d = 2$, conjectural
at $d \geq 3$'') is subsumed by the supertrace formulation
(``$\kappa_{\mathrm{ch}} = \sum_q (-1)^q h^{0,q}$ unconditionally''),
which reduces to the previous $d = 2$ statement and gives the correct
odd-$d$ vanishing. The only genuinely conjectural component is the
identification of $\kappa_{\mathrm{ch}}(\cA_X)$ for
non-product, non-toric, non-local CY$_3$ (Remark
\ref{rem:cy-strat-open-frontier}): in every known family,
$\Xi(X) = \kappa_{\mathrm{ch}}(\cA_X)$ holds, but the verification
requires explicit construction of $\Phi_3(D^b(\Coh(X)))$ on
cases beyond the quintic and the Kummer orbifolds. The conifold
(Corollary~\ref{cor:conifold-non-local-surface}) is the sole
genuinely open case where the direct McKay computation gives a value
($\kappa_{\mathrm{ch}} = 1$) not predicted by $\Xi$, reflecting its
non-compactness.
\end{remark}

\section{Independent verification protocol entries}
\label{sec:cy-strat-independent-verification}

The engine \texttt{compute/lib/cy\_d\_kappa\_d3.py} and the test
file \texttt{compute/tests/test\_cy\_d\_kappa\_stratification.py} install
three disjoint verification paths for Theorem
\ref{thm:kappa-hodge-supertrace-identification} and
Theorem~\ref{thm:kappa-stratification-by-d}:

\begin{enumerate}
\item \textbf{Yau 1977 Calabi conjecture} (path (a)). The Ricci-flat
K\"ahler metric on a CY manifold supplies the $\omega_X \simeq \cO_X$
trivialization used by the HKR--Caldararu side of the computation. The
existence of such a metric is independent of HKR, the chiral bar
complex, or the Vol~I shadow tower; it supplies the $H^q(X, \cO_X)$
column as a Laplacian spectrum.
\item \textbf{Huybrechts \emph{Lectures on K3 Surfaces} 2005} (path
(b)). The K3 Hodge diamond and $\chi(\cO_{K3}) = 2$ from Chapter~1.
Completely disjoint from HKR: the K3 computation is carried out by
explicit Dolbeault cohomology and $c_2(K3) = 24$. This gives the
base of the stratification at $d = 2$.
\item \textbf{Gross--Huybrechts--Joyce \emph{Calabi--Yau Manifolds and
Related Geometries} 2003} (path (c)). The survey chapter on
Bogomolov decomposition and CY$_d$ structure at $d \geq 3$ provides
the column $h^{0,\bullet}$ for the quintic, $\mathrm{K3} \times E$,
$E^3$, and the compact CY$_4$ sextic. Independent of both HKR and
K3-specific arguments.
\end{enumerate}

The three disjoint paths converge on the supertrace values recorded in
Section~\ref{sec:cy-strat-stratification}. See
Section~\ref{sec:cy-strat-stratification} and the test file for the
explicit decorations.

\section{Summary table and closure of Conjecture
\ref{conj:cy-kappa-identification}}
\label{sec:cy-strat-summary}

\begin{center}
\begin{tabular}{lcccc}
\toprule
$X$ & $d$ & $\Xi(X) = \kappa_{\mathrm{ch}}$ & $\chi(\cO_X)$ & match? \\
\midrule
$E$ & $1$ & $0$ & $0$ & yes (cancellation) \\
K3 & $2$ & $2$ & $2$ & \textbf{yes (honest)} \\
$E \times E$ & $2$ & $0$ & $0$ & yes (cancellation) \\
bielliptic & $2$ & $0$ & $0$ & yes (cancellation) \\
quintic $X_5$ & $3$ & $0$ & $0$ & yes (cancellation) \\
K3 $\times E$ & $3$ & $0$ & $0$ & yes (cancellation) \\
$E^3$ & $3$ & $0$ & $0$ & yes (cancellation) \\
local $\P^2$ & $3$ (nc) & $3/2$ & n/a & n/a \\
conifold & $3$ (nc) & $1$ & n/a & n/a \\
sextic $X_6 \subset \P^5$ & $4$ & $2$ & $2$ & yes (honest) \\
generic $CY_5$ & $5$ & $0$ & $0$ & yes (cancellation) \\
\bottomrule
\end{tabular}
\end{center}

Conjecture~\ref{conj:cy-kappa-identification}, as originally stated
(``$\kappa_{\mathrm{ch}} = \chi(\cO_X)$ for all compact CY''), is closed as
follows: the identification holds numerically on the diagonal
$(d, h^{0,\bullet})$ satisfying $h^{0,q} = 0$ for all odd $q$
(equivalently: even $d$ with vanishing Hodge $(1,0), (3,0), \ldots$
columns), and it holds trivially by Serre cancellation on
$(h^{0,q} = h^{0,d-q})$-symmetric profiles (compact CY$_d$ with $d$
odd). The correct \emph{mechanistic} identification is the
supertrace form of
Theorem~\ref{thm:kappa-hodge-supertrace-identification}: $\kappa_{\mathrm{ch}}
(\cA_X) = \sum_q (-1)^q h^{0,q}(X)$. For compact CY of odd dimension,
this is $0$. For compact CY$_d$ with $h^{1,0} = 0$, it equals
$\chi(\cO_X)$. For non-compact toric CY$_d$, it generalizes to a
compactly-supported Dolbeault supertrace and recovers
$\frac{1}{2} \chi_{\mathrm{top}}(\mathrm{base})$ for local surfaces
(Theorem~\ref{thm:local-p2-shadow}); the conifold is the sole
remaining open non-product, non-local case and is handled by direct
McKay quiver cohomology.

The conjecture, rewritten as
\[
 \kappa_{\mathrm{ch}}(\cA_X) \;=\; \sum_{q = 0}^{d} (-1)^q\, h^{0,q}(X)
 \qquad\text{(Conjecture~\ref{conj:cy-kappa-identification}
 in stratified form)},
\]
is promoted to Theorem~\ref{thm:kappa-hodge-supertrace-identification}.
The original \emph{form} of the conjecture (with $\chi(\cO_X)$ on the
right) survives as the $d = 2$, $h^{1,0} = 0$ corollary. The cases
$d \in \{1, 3, 5\}$ record the Serre-enforced vanishing
(Lemma~\ref{lem:chi-O-vanishes-odd-d}). The programme therefore
inscribes the stratification theorem as the unconditional form of
CY-D, with Conjecture~\ref{conj:cy-kappa-identification} closed.
